\section{Methodology}
\label{sec:methodology}

\begin{figure}[t]
    \centering
    \resizebox{0.95\linewidth}{!}{%
    \begin{tabular}{@{}p{0.18\linewidth}p{0.74\linewidth}@{}}
        \toprule
        \textbf{Step} & \textbf{Implementation} \\
        \midrule
        Prompt design & 54 safe late-night cases crossing profile, context length, urgency, and paraphrase. \\
        Generation & 108 real generations from \gptfivemini and \gptfourmini through the \responsesapi. \\
        Labeling & Keyword rules plus an LLM judge separate \nonudge, \softnudge, \harddisengagement, refusal, and sleep-topic answers. \\
        Analysis & Rates, bootstrap intervals, Fisher exact tests with Holm correction, and qualitative error analysis. \\
        \bottomrule
    \end{tabular}
    }
    \caption{Overview of \methodname. The workflow keeps final tasks safe and concrete while varying the user and interaction context, then measures whether the assistant adds a sleep or rest nudge and whether that nudge replaces the requested help.}
    \label{fig:method}
\end{figure}


\para{Target behavior.}
We study assistant-initiated \bedtimenudges: responses that recommend sleep, rest, winding down, calling it a night, or taking a break when the user has asked for safe help. The primary label counts both soft bedtime nudges and hard bedtime disengagements. We separately track whether the assistant answered the request, whether the nudge replaced the answer, and whether the response sounded anxious or overprotective.

\para{Design goal.}
The benchmark tests whether late-night and prolonged-conversation cues trigger a broad response pattern or whether nudges depend on perceived user state. We therefore keep the final tasks safe and concrete, include the same late-night cue in every final user message, and vary profile, context length, urgency, paraphrase, and model.

\begin{table}[t]
    \centering
    \resizebox{\textwidth}{!}{%
    \begin{tabular}{@{}lrl@{}}
        \toprule
        \textbf{Resource} & \textbf{Local rows} & \textbf{Use in this study} \\
        \midrule
        \xstest & 450 & Refusal and exaggerated-safety framing \\
        \donotanswer & 939 & Safety/refusal context and action-label framing \\
        \orbench Hard 1K & 1,319 & Over-refusal context \\
        \orbench Toxic & 655 & Toxic contrast context \\
        \personamem benchmark text & 5,000 & User-profile and long-context prompt inspiration \\
        \wildchat sleep matches & 25 & Observational sleep-phrase audit \\
        \bottomrule
    \end{tabular}
    }
    \caption{Local resources used to design and contextualize the experiment. The causal test uses the synthetic factorial benchmark; the \wildchat sample is an observational audit of phrase-matched conversations.}
    \label{tab:datasets}
\end{table}

\begin{table}[t]
    \centering
    \begin{tabular}{@{}ll@{}}
        \toprule
        \textbf{Factor} & \textbf{Levels} \\
        \midrule
        User profile & neutral, fatigued, dependent \\
        Context length & short (0 prior pairs), medium (4), long (12) \\
        Task urgency & low, urgent \\
        Paraphrase & 3 task variants per urgency level \\
        Model & \gptfivemini, \gptfourmini \\
        \bottomrule
    \end{tabular}
    \caption{Factorial design. The 3 profile $\times$ 3 context $\times$ 2 urgency $\times$ 3 paraphrase design yields 54 prompt cases, each run once on each of two models.}
    \label{tab:design}
\end{table}


\para{Synthetic prompt construction.}
The synthetic benchmark contains 54 prompt cases, each run on two generation models for 108 total generations. The profile factor has three levels. The neutral profile describes a calm, task-focused adult who does not need wellbeing advice unless directly relevant. The fatigued profile describes a tired adult trying to finish safe tasks before stopping for the night. The dependent profile includes vulnerable attachment cues, such as feeling less alone in the chat and feeling uneasy about stopping, but excludes self-harm content. Context length has three levels: short with no prior task pairs, medium with four prior task pairs, and long with twelve prior task pairs. Task urgency has two levels, low and urgent, with three paraphrases per urgency level.

\para{Models and generation settings.}
We used two OpenAI generation models: \gptfivemini (\texttt{gpt-5.4-mini-2026-03-17}) and \gptfourmini (\texttt{gpt-4.1-mini-2025-04-14}), with 54 calls per model. We used the \responsesapi with a 420-token generation cap. The judge model was \gptfourmini with a 300-token cap. Temperature was set to 0 where supported. The random seed was 42 for prompt ordering, bootstrapping, and analysis.

\para{Labeling.}
We label each response in two ways. The keyword baseline searches for explicit sleep and rest phrases, including ``go to sleep,'' ``go to bed,'' ``get some rest,'' ``rest well,'' ``wind down,'' ``call it a night,'' and ``take a break.'' The LLM judge assigns one category from \nonudge, \softnudge, \harddisengagement, ordinary sleep-topic answer, refusal, or crisis escalation. It also records whether the assistant answered the request and rates anxious or overprotective tone from 0 to 3. The primary metric is the assistant-initiated bedtime-nudge rate, defined as the fraction of responses labeled as either \softnudge or \harddisengagement.

\para{Statistical analysis.}
The preregistered plan used logistic regression when event counts permitted, otherwise contingency tests. Because there were only two positive events and 106 negative events, logistic regression was skipped as unstable. We report Fisher exact tests for the planned comparisons and apply Holm correction across the main family of tests. We also report bootstrap 95\% confidence intervals for rate estimates, following the cached analysis outputs.

\para{WildChat audit.}
The observational audit uses a bounded \wildchat sample collected by sleep-phrase matching. The local sample contains 25 conversations selected from 14,782 scanned conversations. We label these conversations with the same judge taxonomy, but we do not use them for causal inference because the sample is selected on phrase matches and includes ordinary sleep-related requests.
